\subsection{Análisis de tráfico de los seis métodos trace}\label{sec:trace-patrones}

Para responder la 3ra pregunta, es decir, qué patrones permiten reconocer las sondas en una captura, revisé los paquetes que cada método emite realmente. Las 6 fichas siguientes usan Valparaíso como ejemplo homogéneo y conservan las comprobaciones de los 3 destinos en la evidencia digital. Todas siguen el mismo esquema, de una ejecución de entrenamiento aprendo la regla de cada campo (como constante, incremento, valor ligado a otro campo), predigo las primeras solicitudes de otra ronda, y las comparo con lo observado. Los valores hexadecimales son bytes de aplicación, los encabezados TCP, UDP e ICMP se analizan por separado.


\Needspace{9\baselineskip}
\subsubsection{scamper / trace / udp}

Tomo Valparaíso como ejemplo. Entrenamiento:  scamper\_\allowbreak{}20261026T204619\_\allowbreak{}270c47c3, contraste con las rondas 2 y 3. Las reglas completas de los tres destinos están en el archivo de predicciones enlazado en el anexo.

\begingroup
\fontsize{9}{11}\selectfont
\addtocounter{table}{-1}
\begin{longtable}{p{3.1cm}p{5.8cm}p{5.1cm}}
\toprule
\textbf{Campo} & \textbf{Regla en ronda 1} & \textbf{Contraste posterior} \\ \midrule
\endfirsthead
\toprule
\textbf{Campo} & \textbf{Regla en ronda 1} & \textbf{Contraste posterior} \\ \midrule
\endhead
ip.ttl & incremento 1 por paquete & r2: cumple; r3: cumple \\
ip.id & incremento 1 por paquete & r2: cumple; r3: cumple \\
ip.len & constante 44 & r2: cumple; r3: cumple \\
ip.proto & constante 17 & r2: cumple; r3: cumple \\
udp.srcport & constante 32769 & r2: cumple; r3: cumple \\
udp.dstport & incremento 1 por paquete & r2: cumple; r3: cumple \\
udp.length & constante 24 & r2: cumple; r3: cumple \\
udp.checksum & incremento -1 por paquete & r2: cumple; r3: cumple \\
\bottomrule
\end{longtable}
\endgroup
Payload: [16] bytes. Hexadecimal inicial (offset 0):

\begin{quote}\ttfamily\footnotesize\raggedright 0000 0000 0000 0000 0000 0000 0000 0000\end{quote}

%La suma de payload\_\allowbreak{}word0\_\allowbreak{}BE, udp.checksum, udp.srcport, udp.dstport módulo 65535 se mantuvo constante. La base se aprende del primer paquete de cada ejecución; no se presume universal. Es compatible con compensación de checksum, sin probar la intención del emisor.

%Predicción concreta en ronda 2, campo udp.dstport: la primera solicitud tiene valor 33435. Con la regla aprendida y esa base cuando corresponde, se predicen las solicitudes 2 a 4 antes de contrastarlas con sus valores observados.

Al sumar payload\_word0\_BE, udp.checksum, udp.srcport y udp.dstport módulo 65535, el resultado se mantiene constante durante la ejecución. Ese valor base se aprende del primer paquete de cada ejecución y no lo doy por universal. Es lo que se esperaría si el emisor compensara el checksum con el payload, aunque con esto no se puede probar que esa fuera su intención. 
Predicción en la ronda 2, campo udp.dstport: la primera solicitud tiene valor 33435. Aplicando la regla aprendida (y esa base, cuando corresponde) predije los valores de las solicitudes 2 a 4 antes de compararlos con lo que realmente se observó.


\begingroup
\fontsize{9}{11}\selectfont
\addtocounter{table}{-1}
\begin{longtable}{p{3.2cm}p{5.1cm}p{5.1cm}}
\toprule
\textbf{Solicitud / trama} & \textbf{Predicción} & \textbf{Observación} \\ \midrule
\endfirsthead
\toprule
\textbf{Solicitud / trama} & \textbf{Predicción} & \textbf{Observación} \\ \midrule
\endhead
2 / 7 & 33436 & 33436 \\
3 / 9 & 33437 & 33437 \\
4 / 11 & 33438 & 33438 \\
\bottomrule
\end{longtable}
\endgroup
Alcance: 9 solicitudes del ejemplo inicial. Verificar un campo o un byte no es igual que una medición independiente, y a los campos que varían les dejo sus restricciones en lugar de inventarles un valor. 


\Needspace{9\baselineskip}
\subsubsection{scamper / trace / icmp}

%Ejemplo: Valparaíso. Entrenamiento: scamper\_\allowbreak{}20261026T204623\_\allowbreak{}0cf6d8d5; contraste: rondas 2 y 3. Las reglas completas de los tres destinos están en el archivo de predicciones enlazado en el anexo.

Valparaíso como ejemplo. Entrenamiento: scamper\_20261026T204623\_0cf6d8d5; contraste con las rondas 2 y 3. 


\begingroup
\fontsize{9}{11}\selectfont
\addtocounter{table}{-1}
\begin{longtable}{p{3.1cm}p{5.8cm}p{5.1cm}}
\toprule
\textbf{Campo} & \textbf{Regla en ronda 1} & \textbf{Contraste posterior} \\ \midrule
\endfirsthead
\toprule
\textbf{Campo} & \textbf{Regla en ronda 1} & \textbf{Contraste posterior} \\ \midrule
\endhead
ip.ttl & incremento 1 por paquete & r2: cumple; r3: cumple \\
ip.id & incremento 1 por paquete & r2: cumple; r3: cumple \\
ip.len & constante 44 & r2: cumple; r3: cumple \\
ip.proto & constante 1 & r2: cumple; r3: cumple \\
icmp.type & constante 8 & r2: cumple; r3: cumple \\
icmp.code & constante 0 & r2: cumple; r3: cumple \\
icmp.ident & constante 32769 & r2: cumple; r3: cumple \\
icmp.seq & incremento 1 por paquete & r2: cumple; r3: cumple \\
icmp.checksum & incremento -1 por paquete & r2: cumple; r3: cumple \\
\bottomrule
\end{longtable}
\endgroup
Payload: [16] bytes. Hexadecimal inicial (offset 0):

%La suma de payload\_\allowbreak{}word0\_\allowbreak{}BE, icmp.checksum, icmp.ident, icmp.seq módulo 65535 se mantuvo constante. La base se aprende del primer paquete de cada ejecución; no se presume universal. Es compatible con compensación de checksum, sin probar la intención del emisor.

Al sumar payload\_word0\_BE, icmp.checksum, icmp.ident y icmp.seq módulo 65535, el resultado se mantiene constante durante la ejecución. Ese valor base se aprende del primer paquete de cada ejecución y no lo doy por universal. Es lo que se esperaría si el emisor compensara el checksum con el payload, aunque con esto no puedo probar que esa fuera su intención.

Predicción concreta en ronda 2, campo icmp.seq: la primera solicitud tiene valor 0. Aplicando la regla aprendida (y esa base cuando corresponde), se predicen las solicitudes 2 a 4 antes de contrastarlas con sus valores observados.

\begingroup
\fontsize{9}{11}\selectfont
\addtocounter{table}{-1}
\begin{longtable}{p{3.2cm}p{5.1cm}p{5.1cm}}
\toprule
\textbf{Solicitud / trama} & \textbf{Predicción} & \textbf{Observación} \\ \midrule
\endfirsthead
\toprule
\textbf{Solicitud / trama} & \textbf{Predicción} & \textbf{Observación} \\ \midrule
\endhead
2 / 5 & 1 & 1 \\
3 / 7 & 2 & 2 \\
4 / 9 & 3 & 3 \\
\bottomrule
\end{longtable}
\endgroup
Alcance: 9 solicitudes del ejemplo inicial. Verificar un campo o un byte no es lo mismo que una medición independiente, y a los campos que varían les dejo sus restricciones en lugar de inventarles un valor.

\Needspace{9\baselineskip}
\subsubsection{scamper / trace / udp-paris}
Ejemplo: Valparaíso. Entrenamiento: scamper\_\allowbreak{}20261026T204627\_\allowbreak{}61f04669; contraste: rondas 2 y 3. 

\begingroup
\fontsize{9}{11}\selectfont
\addtocounter{table}{-1}
\begin{longtable}{p{3.1cm}p{5.8cm}p{5.1cm}}
\toprule
\textbf{Campo} & \textbf{Regla en ronda 1} & \textbf{Contraste posterior} \\ \midrule
\endfirsthead
\toprule
\textbf{Campo} & \textbf{Regla en ronda 1} & \textbf{Contraste posterior} \\ \midrule
\endhead
ip.ttl & incremento 1 por paquete & r2: cumple; r3: cumple \\
ip.id & incremento 1 por paquete & r2: cumple; r3: cumple \\
ip.len & constante 44 & r2: cumple; r3: cumple \\
ip.proto & constante 17 & r2: cumple; r3: cumple \\
udp.srcport & constante 32769 & r2: cumple; r3: cumple \\
udp.dstport & constante 33435 & r2: cumple; r3: cumple \\
udp.length & constante 24 & r2: cumple; r3: cumple \\
udp.checksum & incremento 1 por paquete & r2: cumple; r3: cumple \\
\bottomrule
\end{longtable}
\endgroup
Payload: [16] bytes. Hexadecimal inicial (offset 0):

\begin{quote}\ttfamily\footnotesize\raggedright 66b1 0000 0000 0000 0000 0000 0000 0000\end{quote}
%La suma de payload\_\allowbreak{}word0\_\allowbreak{}BE, udp.checksum, udp.srcport, udp.dstport módulo 65535 se mantuvo constante. La base se aprende del primer paquete de cada ejecución; no se presume universal. Es compatible con compensación de checksum, sin probar la intención del emisor.

Al sumar payload\_word0\_BE, udp.checksum, udp.srcport y udp.dstport módulo 65535, el resultado se mantiene constante durante la ejecución. Ese valor base se aprende del primer paquete de cada ejecución y no lo doy por universal. Es lo que se esperaría si el emisor compensara el checksum con el payload, aunque con esto no puedo probar que esa fuera su intención.

Predicción en la ronda 2, campo ip.ttl: la primera solicitud tiene valor 1. Aplicando la regla aprendida (y esa base, cuando corresponde) predije los valores de las solicitudes 2 a 4 antes de compararlos con lo que realmente se observó.

\begingroup
\fontsize{9}{11}\selectfont
\addtocounter{table}{-1}
\begin{longtable}{p{3.2cm}p{5.1cm}p{5.1cm}}
\toprule
\textbf{Solicitud / trama} & \textbf{Predicción} & \textbf{Observación} \\ \midrule
\endfirsthead
\toprule
\textbf{Solicitud / trama} & \textbf{Predicción} & \textbf{Observación} \\ \midrule
\endhead
2 / 10 & 2 & 2 \\
3 / 13 & 3 & 3 \\
4 / 15 & 4 & 4 \\
\bottomrule
\end{longtable}
\endgroup

Alcance: 9 solicitudes del ejemplo inicial. Verificar un campo o un byte no es lo mismo que una medición independiente, y a los campos que varían les dejo sus restricciones en lugar de inventarles un valor.

\Needspace{9\baselineskip}
\subsubsection{scamper / trace / icmp-paris}
Ejemplo: Valparaíso. Entrenamiento: scamper\_\allowbreak{}20261026T204633\_\allowbreak{}256949f1; contraste: rondas 2 y 3. 

\begingroup
\fontsize{9}{11}\selectfont
\addtocounter{table}{-1}
\begin{longtable}{p{3.1cm}p{5.8cm}p{5.1cm}}
\toprule
\textbf{Campo} & \textbf{Regla en ronda 1} & \textbf{Contraste posterior} \\ \midrule
\endfirsthead
\toprule
\textbf{Campo} & \textbf{Regla en ronda 1} & \textbf{Contraste posterior} \\ \midrule
\endhead
ip.ttl & incremento 1 por paquete & r2: cumple; r3: cumple \\
ip.id & incremento 1 por paquete & r2: cumple; r3: cumple \\
ip.len & constante 44 & r2: cumple; r3: cumple \\
ip.proto & constante 1 & r2: cumple; r3: cumple \\
icmp.type & constante 8 & r2: cumple; r3: cumple \\
icmp.code & constante 0 & r2: cumple; r3: cumple \\
icmp.ident & constante 32769 & r2: cumple; r3: cumple \\
icmp.seq & incremento 1 por paquete & r2: cumple; r3: cumple \\
icmp.checksum & constante 0x8001 & r2: cumple; r3: cumple \\
\bottomrule
\end{longtable}
\endgroup
Payload: [16] bytes. Hexadecimal inicial (offset 0):

\begin{quote}\ttfamily\footnotesize\raggedright f7fc 0000 0000 0000 0000 0000 0000 0000\end{quote}
La suma de payload\_\allowbreak{}word0\_\allowbreak{}BE, icmp.checksum, icmp.ident, icmp.seq módulo 65535 se mantuvo constante. Ese valor base se aprende del primer paquete de cada ejecución, evitando presumir universal. Es compatible con compensación de checksum, sin probar la intención del emisor.

Predicción en la ronda 2, campo icmp.seq: la primera solicitud tiene valor 0. En base a la regla aprendida, predije los valores de las solicitudes 2 a 4 antes de compararlos con lo que realmente se observó.

\begingroup
\fontsize{9}{11}\selectfont
\addtocounter{table}{-1}
\begin{longtable}{p{3.2cm}p{5.1cm}p{5.1cm}}
\toprule
\textbf{Solicitud / trama} & \textbf{Predicción} & \textbf{Observación} \\ \midrule
\endfirsthead
\toprule
\textbf{Solicitud / trama} & \textbf{Predicción} & \textbf{Observación} \\ \midrule
\endhead
2 / 6 & 1 & 1 \\
3 / 10 & 2 & 2 \\
4 / 12 & 3 & 3 \\
\bottomrule
\end{longtable}
\endgroup
Alcance: 9 solicitudes en el ejemplo inicial. Una verificación de campo o byte no equivale a una medición independiente; los campos variables conservan sus restricciones y no reciben valores inventados.

\Needspace{9\baselineskip}
\subsubsection{scamper / trace / tcp}
Ejemplo: Valparaíso. Entrenamiento: scamper\_\allowbreak{}20261026T204638\_\allowbreak{}04ca1c5e; contraste: rondas 2 y 3. 

\begingroup
\fontsize{9}{11}\selectfont
\addtocounter{table}{-1}
\begin{longtable}{p{3.1cm}p{5.8cm}p{5.1cm}}
\toprule
\textbf{Campo} & \textbf{Regla en ronda 1} & \textbf{Contraste posterior} \\ \midrule
\endfirsthead
\toprule
\textbf{Campo} & \textbf{Regla en ronda 1} & \textbf{Contraste posterior} \\ \midrule
\endhead
ip.ttl & incremento 1 cada 2 paquetes & r2: cumple; r3: cumple \\
ip.id & incremento 1 por paquete & r2: cumple; r3: cumple \\
ip.len & constante 40 & r2: cumple; r3: cumple \\
ip.proto & constante 6 & r2: cumple; r3: cumple \\
tcp.srcport & constante 32769 & r2: cumple; r3: cumple \\
tcp.dstport & constante 80 & r2: cumple; r3: cumple \\
tcp.flags & constante 0x0002 & r2: cumple; r3: cumple \\
tcp.seq\_\allowbreak{}raw & constante 0 & r2: cumple; r3: cumple \\
tcp.ack\_\allowbreak{}raw & constante 0 & r2: cumple; r3: cumple \\
tcp.len & constante 0 & r2: cumple; r3: cumple \\
tcp.checksum & constante 0x9922 & r2: cumple; r3: cumple \\
\bottomrule
\end{longtable}
\endgroup
Payload: [0] bytes. No hay datos de aplicación en las sondas TCP observadas.

Predicción en la ronda 2, campo ip.ttl: la primera solicitud tiene valor 1.

\begingroup
\fontsize{9}{11}\selectfont
\addtocounter{table}{-1}
\begin{longtable}{p{3.2cm}p{5.1cm}p{5.1cm}}
\toprule
\textbf{Solicitud / trama} & \textbf{Predicción} & \textbf{Observación} \\ \midrule
\endfirsthead
\toprule
\textbf{Solicitud / trama} & \textbf{Predicción} & \textbf{Observación} \\ \midrule
\endhead
2 / 30 & 1 & 1 \\
3 / 41 & 2 & 2 \\
4 / 54 & 2 & 2 \\
\bottomrule
\end{longtable}
\endgroup
Alcance: 10 solicitudes en el ejemplo inicial. Una verificación de campo o byte no equivale a una medición independiente, cuyos campos variables conservan sus restricciones y no reciben valores inventados.

\Needspace{9\baselineskip}
\subsubsection{scamper / trace / tcp-ack}
Ejemplo: Valparaíso. Entrenamiento: scamper\_\allowbreak{}20261026T204732\_\allowbreak{}c33e1983; contraste: rondas 2 y 3. 

\begingroup
\fontsize{9}{11}\selectfont
\addtocounter{table}{-1}
\begin{longtable}{p{3.1cm}p{5.8cm}p{5.1cm}}
\toprule
\textbf{Campo} & \textbf{Regla en ronda 1} & \textbf{Contraste posterior} \\ \midrule
\endfirsthead
\toprule
\textbf{Campo} & \textbf{Regla en ronda 1} & \textbf{Contraste posterior} \\ \midrule
\endhead
ip.ttl & incremento 1 cada 2 paquetes & r2: cumple; r3: cumple \\
ip.id & incremento 1 por paquete & r2: cumple; r3: cumple \\
ip.len & constante 40 & r2: cumple; r3: cumple \\
ip.proto & constante 6 & r2: cumple; r3: cumple \\
tcp.srcport & constante 32769 & r2: cumple; r3: cumple \\
tcp.dstport & constante 80 & r2: cumple; r3: cumple \\
tcp.flags & constante 0x0010 & r2: cumple; r3: cumple \\
tcp.seq\_\allowbreak{}raw & constante 0 & r2: cumple; r3: cumple \\
tcp.ack\_\allowbreak{}raw & constante 0 & r2: cumple; r3: cumple \\
tcp.len & constante 0 & r2: cumple; r3: cumple \\
tcp.checksum & constante 0x9914 & r2: cumple; r3: cumple \\
\bottomrule
\end{longtable}
\endgroup
Payload: [0] bytes. No hay datos de aplicación en las sondas TCP observadas.

Predicción concreta en ronda 2, campo ip.ttl: la primera solicitud tiene valor 1. 

\begingroup
\fontsize{9}{11}\selectfont
\addtocounter{table}{-1}
\begin{longtable}{p{3.2cm}p{5.1cm}p{5.1cm}}
\toprule
\textbf{Solicitud / trama} & \textbf{Predicción} & \textbf{Observación} \\ \midrule
\endfirsthead
\toprule
\textbf{Solicitud / trama} & \textbf{Predicción} & \textbf{Observación} \\ \midrule
\endhead
2 / 4 & 1 & 1 \\
3 / 8 & 2 & 2 \\
4 / 10 & 2 & 2 \\
\bottomrule
\end{longtable}
\endgroup
Alcance: 10 solicitudes en el ejemplo inicial. 

\subsubsection{Lectura conjunta de los seis métodos}

En conjunto, los 6 métodos se distinguen por pocos rasgos. En udp e icmp clásicos el payload son 16 bytes en cero y lo que cambia entre solicitudes es el puerto de destino (que parte en 33435) o el icmp.seq (que parte en 0), de modo que cada sonda es un flujo distinto para un balanceador. En udp-paris e icmp-paris los dos primeros bytes del payload dejan de ser 0 (66b1 y f7fc en los ejemplos) y el resto sigue en cero, y lo que se conserva es la suma con el checksum y los campos de flujo (es justo lo que se espera de las variantes Paris, que intentan mantener constante lo que usan los balanceadores para elegir el camino). Con tcp y tcp-ack no hay datos de aplicación y las sondas se reconocen por los flags (SYN 0x0002 y ACK 0x0010) y por un TTL que parte en 1.

Conviene matizar la lectura de la suma módulo 65535. Cuando solo cambian esos campos, la aritmética del checksum hace que la suma se conserve por sí sola, por lo que la regularidad observada es coherente con una compensación deliberada pero no la prueba. Lo que sí queda establecido es que las reglas aprendidas en una ronda predijeron los campos de las solicitudes siguientes, permitiendo reconocer estas sondas en una captura sin depender de la salida de la herramienta.



